\chapter{光场也改变时，必须追踪电子与光的联合态}
\label{chap:feqo-quantum-exchange}

经典光场的一个复数 \(\beta\) 能给出电子边带，却不能回答交换之后光子数怎样变化。若一束光只含少数光子，吸收一个光子会明显改变光场；电子与光的末态必须一起求出。这里沿用上一章的电子边带标号 \(\ell\)，并使用已在场量子化中定义的单模数态 \(|n\rangle_\gamma\)。

电子与光的联合态属于张量积空间 \(\mathcal H_e\otimes\mathcal H_\gamma\)。在积基 \(|\ell,n\rangle=|\ell\rangle_e\otimes|n\rangle_\gamma\) 中，任意归一化纯态可写为 \(|\Psi\rangle=\sum_{\ell,n}C_{\ell n}|\ell,n\rangle\)，且 \(\sum|C_{\ell n}|^2=1\)。同一电子态若能与唯一光态相乘写出，称可分；若做不到，称纠缠。接下来用一次交换实际计算这个差别。

实验通常只读电子。要从联合态得到电子的所有测量概率，需要一个能保留电子相干项、又自动对未读出的光结果求和的对象。若制备本身还混合了几种可能的纯态，先用密度算符记录这种经典不确定性。

\begin{definition}[密度算符与偏迹]
制备 \(|\psi_k\rangle\) 的概率为 \(p_k\geq0\)、\(\sum_kp_k=1\) 时，密度算符为 \(\rho=\sum_kp_k|\psi_k\rangle\langle\psi_k|\)。它正、迹为一，并以 \(\operatorname{Tr}(\rho A)\) 给可观测量 \(A\) 的期望。联合态只测电子时，电子约化态 \(\rho_e=\operatorname{Tr}_\gamma\rho_{e\gamma}\) 被定义为满足
\(\operatorname{Tr}_e(\rho_e A)=\operatorname{Tr}_{e\gamma}[\rho_{e\gamma}(A\otimes I)]\) 的唯一算符。
\end{definition}

取任意光子数正交基，将 \(|\Psi\rangle=\sum C_{\ell n}|\ell,n\rangle\) 代入定义，立即得到
\begin{equation}
(\rho_e)_{\ell\ell'}=\sum_n C_{\ell n}C_{\ell'n}^*.
\label{eq:feqo-partial-trace-components}
\end{equation}
固定末态光子数 \(n\) 后振幅才相乘；不同 \(n\) 之间求的是概率项之和，不是把所有振幅先加起来。这一行是后面比较不同光统计的共同计算规则。

\section{一个量子交换会改变两边的量子数}
\label{chap:feqo-quantum-pinem}

令电子边带数算符 \(L_e|\ell\rangle=\ell|\ell\rangle\)，平移算符 \(B|\ell\rangle=|\ell-1\rangle\)、\(B^\dagger|\ell\rangle=|\ell+1\rangle\)。在双向无限的理想电子梯上 \(B^\dagger B=BB^\dagger=I\)；它不是满足 \([B,B^\dagger]=1\) 的玻色湮灭算符。光模仍由 \(a|n\rangle=\sqrt n|n-1\rangle\) 描述。

在低反冲、单模、旋波条件下，单次能量交换的 哈密顿量 可写为
\begin{equation}
H_I(t)=\hbar g(t)B^\dagger a+\hbar g(t)^*B a^\dagger.
\label{eq:feqo-quantum-pinem-hamiltonian}
\end{equation}
\(g\) 单位为 \(\si{\per\second}\)。第一项提高电子边带标号并湮灭一个光子，第二项反向进行。令 \(N_\gamma=a^\dagger a\)。由 \([L_e,B^\dagger]=B^\dagger\)、\([N_\gamma,a]=-a\)，直接相加得到
\begin{equation}
[H_I,N_\gamma+L_e]=0.
\label{eq:feqo-exchange-conservation}
\end{equation}
守恒的是 \(N_\gamma+L_e\)，不是 \(N_\gamma+B^\dagger B\)，后者在无限梯上仅等于 \(N_\gamma+I\)。这个区别决定哪些联合态可以相互连接。

将 \(|\Psi(t)\rangle=\sum C_{\ell n}(t)|\ell,n\rangle\) 代入薛定谔方程，分别使两项作用于源态，得到
\begin{equation}
\ii\dot C_{\ell n}=g(t)\sqrt{n+1}\,C_{\ell-1,n+1}
+g(t)^*\sqrt n\,C_{\ell+1,n-1}.
\label{eq:feqo-quantum-pinem-recursion}
\end{equation}
右边两项都满足 \((\ell-1)+(n+1)=\ell+n=(\ell+1)+(n-1)\)。截断数值计算时要在光子数上保留足够大的上界，并检查边界人口；简单把越界振幅扔掉会破坏幺正性。

最小手算例子需要明确边界。若一个物理选择规则或强失谐只保留 \(|0,1\rangle\) 与 \(|1,0\rangle\)，用投影 \(P\) 定义受限模型 \(PH_IP\)，常数实 \(g\) 给二维矩阵 \(\hbar g\sigma_x\)。从 \(|0,1\rangle\) 出发，
\[
|\Psi(t)\rangle=\cos(gt)|0,1\rangle-\ii\sin(gt)|1,0\rangle.
\]
由式~\eqref{eq:feqo-partial-trace-components}，\(\rho_e=\cos^2(gt)|0\rangle\langle0|+\sin^2(gt)|1\rangle\langle1|\)，纯度为 \(\operatorname{Tr}\rho_e^2=1-\tfrac12\sin^2(2gt)\)。联合态始终是纯态，电子在中间时刻却是混态。若无限电子梯未被物理地截断，\(|-1,2\rangle\) 也与 \(|0,1\rangle\) 相连，上面的二维解就不是完整 哈密顿量 的精确解。

有限截断时还有一个有用的矩阵定理，值得在这里说明其含义而非只报名字。任意有限复矩阵 \(C\) 都可写成 \(C=\sum_r s_r u_rv_r^\dagger\)，其中 \(s_r>0\)，两组列向量 \(u_r,v_r\) 各自正交归一；这叫奇异值分解。证明从正的 厄米 矩阵 \(C^\dagger C\) 出发：取其非零本征值 \(s_r^2\) 与归一本征向量 \(v_r\)，令 \(u_r=Cv_r/s_r\)。于是 \(u_r^\dagger u_s=v_r^\dagger C^\dagger Cv_s/(s_rs_s)=\delta_{rs}\)；在 \(C^\dagger C\) 的零空间上 \(C\) 为零，所以这些非零项已重构整个 \(C\)。

把矩阵分解代回联合态，光侧向量的分量取 \(v_r\) 的复共轭，便得 \(|\Psi\rangle=\sum_r s_r|u_r\rangle_e|\bar v_r\rangle_\gamma\)。矩阵归一化给 \(\sum_rs_r^2=\sum_{\ell,n}|C_{\ell n}|^2=1\)。把它代入偏迹式得到 \(\rho_e=\sum_rs_r^2|u_r\rangle\langle u_r|\)。因此有限截断态可分当且仅当只有一个非零奇异值；\(\operatorname{Tr}\rho_e^2=\sum_rs_r^4\)。奇异值分解只在这个有限矩阵上使用，不能不加说明地把有限维定理搬到无限电子梯。

无限梯的判据其实可直接由偏迹证明，不需要先调用无限维奇异值定理。把矩阵的第 \(n\) 列看成电子向量 \(|c_n\rangle_e=\sum_\ell C_{\ell n}|\ell\rangle\)，则归一化条件为 \(\sum_n\|c_n\|^2=1\)，而式~\eqref{eq:feqo-partial-trace-components} 成为 \(\rho_e=\sum_n|c_n\rangle\langle c_n|\)。因为所有求和项非负，可先对有限部分求迹再取极限，得到
\begin{equation}
\operatorname{Tr}\rho_e^2
=\sum_{n,m}|\langle c_n|c_m\rangle|^2
\leq\sum_{n,m}\|c_n\|^2\|c_m\|^2=1.
\label{eq:feqo-entanglement-column-test}
\end{equation}
不等号逐项使用柯西--施瓦茨不等式。等号成立当且仅当所有非零列向量彼此共线；此时 \(C_{\ell n}=u_\ell v_n\)，联合态正好是 \((\sum_\ell u_\ell|\ell\rangle)\otimes(\sum_n v_n|n\rangle)\)。若至少两列不共线，严格不等号出现，电子约化态的纯度小于一，联合纯态便是纠缠态。这个证明对平方可和的无限矩阵同样成立：有限部分单调收敛，且右边总和为一，不需要对无限矩阵做未经证明的谱分解。

理想无限电子梯还有一个比二维截断更有用的精确解。令 \(A=B^\dagger a\)，则 \(A^\dagger=Ba^\dagger\)。因为电子平移与光算符作用在不同空间，且 \(BB^\dagger=B^\dagger B=I\)，
\[
[A,A^\dagger]=B^\dagger B\,aa^\dagger-BB^\dagger\,a^\dagger a=I.
\]
满足玻色对易关系的是电子与光组合成的 \(A\)，不是电子平移 \(B\)。若 \(g\) 在持续时间 \(T\) 内为常数，令 \(\zeta=-\ii g^*T\)，则式~\eqref{eq:feqo-quantum-pinem-hamiltonian} 的精确传播为
\begin{equation}
U(T)=\exp(\zeta A^\dagger-\zeta^*A)
=e^{-|\zeta|^2/2}e^{\zeta A^\dagger}e^{-\zeta^*A}.
\label{eq:feqo-joint-displacement}
\end{equation}
第二个等号是对易子为数时的 Baker--Campbell--Hausdorff 恒等式。它可直接验证：将两个指数的乘积按级数展开，所有高阶对易子都为零，只剩 \(-|\zeta|^2/2\)；而指数中的算符反厄米，故 \(U^\dagger U=I\)。在初态 \(|0,0\rangle\) 上，\(A|0,0\rangle=0\)，逐项展开右边指数便有
\begin{equation}
U(T)|0,0\rangle=e^{-\mu/2}\sum_{m=0}^\infty
\frac{\zeta^m}{\sqrt{m!}}|-m,m\rangle,
\qquad \mu=|\zeta|^2.
\label{eq:feqo-vacuum-joint-state}
\end{equation}
电子失去 \(m\hbar\omega\)，同一光模多出 \(m\) 个光子。按式~\eqref{eq:feqo-partial-trace-components} 求偏迹，电子人口为 \(P_{-m}=e^{-\mu}\mu^m/m!\)，其它边带人口为零；\(\sum_mP_{-m}=e^{-\mu}e^\mu=1\)。约化纯度是 \(e^{-2\mu}I_0(2\mu)\)，其中 \(I_0(2\mu):=\sum_{m=0}^\infty\mu^{2m}/(m!)^2\)。\(\mu>0\) 时其值小于一，所以这个全阶例子展示了电子与光的纠缠。

这一联合态还能把“交换”二字变成可测的相关。以 \(P_m=e^{-\mu}\mu^m/m!\) 求生成函数 \(F(s)=\sum_mP_ms^m=e^{\mu(s-1)}\)，在 \(s=1\) 求一次和两次导数，便得 \(\langle m\rangle=\mu\)、\(\langle m(m-1)\rangle=\mu^2\)，所以 \(\operatorname{Var}(m)=\mu\)。每次事例中 \(L_e=-N_\gamma=-m\)，故
\[
\langle\Delta E_e\rangle=-\hbar\omega\mu,\qquad
\operatorname{Cov}(L_e,N_\gamma)=-\mu.
\]
光子计数与电子损失量逐次严格反相关，即使只看电子或只看光的边缘分布都会丢掉这条信息。这也是实际检验单模交换模型的一个办法：若同时记录两侧而相关显著偏离上述守恒线，须检查未收集的光模、电子反冲或探测器漏计。

这个位移解依赖四项条件：电子梯在实际占据区可视作双向无限；各交换边的反冲失谐可忽略；只保留一个共振光模；\(g\) 的相位在 \(T\) 内固定。若 \(g(t)\) 仅实幅度变化而相位固定，可把 \(gT\) 换成 \(\int g(t)\dd t\)；若相位变化，时间排序产生额外整体相位。若反冲可分辨，自由失谐不再与相互作用对易，式~\eqref{eq:feqo-joint-displacement} 就不是完整传播。

单模位移把一个共振频率视为唯一可用的光通道。实际结构的模常常连续；这时不能把不同频率的振幅无条件相加。先在长度有限的量子化盒中取离散模 \(a_j\)，\([a_j,a_k^\dagger]=\delta_{jk}\)，每个模具有频率 \(\omega_j\)、动量转移 \(\hbar\mathbf q_j\) 和角频率量纲耦合 \(g_j\)。一次吸收从 \(|\mathbf p;n_j\rangle\) 到 \(|\mathbf p+\hbar\mathbf q_j;n_j-1\rangle\) 的一阶振幅是
\begin{equation}
\mathcal A_j=-\ii g_j\sqrt{n_j}\,
\widetilde w(\Delta_j),\qquad
\Delta_j=\frac{E(\mathbf p+\hbar\mathbf q_j)-E(\mathbf p)}{\hbar}-\omega_j .
\label{eq:feqo-continuum-mode-amplitude}
\end{equation}
\(\widetilde w\) 是式~\eqref{eq:feqo-finite-time-amplitude} 定义的时间窗变换。若初态各模是确定数态，\(j\neq k\) 的末态光场正交，故总一次吸收概率是 \(\sum_j|\mathcal A_j|^2\)；把振幅先求和再取模平方会制造不存在的模式干涉。

令相邻频率间隔为 \(\Delta\omega\)，定义 \(a(\omega_j)=a_j/\sqrt{\Delta\omega}\)、\(g(\omega_j)=g_j/\sqrt{\Delta\omega}\)。当盒变大，\([a(\omega),a^\dagger(\omega')]=\delta(\omega-\omega')\)，并且 \(\sum_jg_ja_j\to\int\dd\omega\,g(\omega)a(\omega)\)。\(a(\omega)\) 的量纲是 \(\si{\second\tothe{1/2}}\)，\(g(\omega)\) 是 \(\si{\per\second\tothe{1/2}}\)；两者连同 \(\dd\omega\) 仍给角频率。对模式族标号 \(\nu\) 求和后，定义频率谱密度 \(J(\omega)=\sum_\nu|g_\nu(\omega)|^2\)，量纲为 \(\si{\per\second}\)。若不同 \(\nu\) 的失谐只取决于 \(\omega\)，
\begin{equation}
P_{\rm abs}^{(1)}=\int_0^\infty\dd\omega\,
n(\omega)J(\omega)|\widetilde w[\Delta(\omega)]|^2.
\label{eq:feqo-continuum-absorption}
\end{equation}
对真空发射，相关因子不是 \(n\) 而是 \(n+1\)。若 \(\Delta(\omega_*)=0\) 的根孤立且 \(\Delta'(\omega_*)\neq0\)，矩形窗的长时极限给 \(P/T\to2\pi\sum_{\omega_*}n(\omega_*)J(\omega_*)/|\Delta'(\omega_*)|\)。这里的导数是从 \(\delta[\Delta(\omega)]\) 换到 \(\delta(\omega-\omega_*)\) 时的 Jacobian；若 \(J,n\) 在宽度 \(1/T\) 内变化剧烈、\(\Delta'=0\) 或总跃迁概率不再小，就不能使用这个速率。

电子探测器通常按能量而非动量分箱，另一个 Jacobian 来自态归一化。固定横向动量和自旋，沿正向一维支有 \(\langle p|p'\rangle=\delta(p-p')\)、\(v(E)=\dd E/\dd p>0\)。定义 \(|E\rangle_E=v(E)^{-1/2}|p(E)\rangle\)，则
\[
{}_E\langle E|E'\rangle_E=
\frac{\delta[p(E)-p(E')]}{\sqrt{v(E)v(E')}}
=\delta(E-E').
\]
因此动量概率密度 \(|f_p(p)|^2\dd p\) 变成 \(|f_p[p(E)]|^2\dd E/v(E)\)。这个速度因子不能在从理论动量态改成实验能谱时遗漏。固定横向动量的限定也必须保留；若收集了不同出射角，还须对横向动量积分并使用相应的完整相空间测度。

\begin{exercise}
检查 \(H_I\) 对 \(|0,1\rangle\) 的完整作用，并解释为什么除 \(|1,0\rangle\) 外还出现 \(|-1,2\rangle\)。再从 \(\delta[\Delta(\omega)]\) 推导连续模速率中的 \(|\Delta'|^{-1}\)。
\end{exercise}


一次交换已经让电子与光可以纠缠；即使平均光子数相同，不同光初态的相位和涨落也不相同。现在把电子测量当作探针，问它分别读取光的哪个统计量。

\section{同样平均光子数，为什么电子相干性仍不同}
\label{chap:feqo-photon-statistics}

把电子初态取为 \(|0\rangle_e\)，光的初态取任意单模密度算符 \(\rho_\gamma\)。在短时间或弱耦合下设 \(\xi=\int g(t)\dd t\)，并要求 \(|\xi|\sqrt{\langle N_\gamma\rangle+1}\ll1\)。式~\eqref{eq:feqo-quantum-pinem-hamiltonian} 的一阶传播为
\[
U=I-\ii\xi B^\dagger a-\ii\xi^*B a^\dagger+O(|\xi|^2).
\]
为求到 \(|\xi|^2\) 阶的人口，还需把二阶返回原态的振幅计入；利用幺正性确定其总和即可。记 \(\bar n=\operatorname{Tr}(\rho_\gamma N_\gamma)\)，则
\begin{equation}
P_{+1}=|\xi|^2\bar n+O(|\xi|^4),\quad
P_{-1}=|\xi|^2(\bar n+1)+O(|\xi|^4),\quad
P_0=1-|\xi|^2(2\bar n+1)+O(|\xi|^4).
\label{eq:feqo-weak-quantum-populations}
\end{equation}
这里 \(a^\dagger a=N_\gamma\) 给吸收权重，\(aa^\dagger=N_\gamma+1\) 给发射权重；额外的 1 是真空涨落。若有初末态失谐，\(\xi\) 应换成相应的有限时间 Fourier 积分，不能只积 \(g\)。

为什么更高边带能够区分统计？在相位固定且忽略反冲的相互作用中，第二次吸收来自传播指数的平方项 \(-\xi^2(B^\dagger a)^2/2\)。初始光子数为 \(n\) 时，\(a^2|n\rangle=\sqrt{n(n-1)}|n-2\rangle\)，于是对任意光初态，领先的双吸收概率为
\begin{equation}
P_{+2}=\frac{|\xi|^4}{4}
\langle N_\gamma(N_\gamma-1)\rangle+O(|\xi|^6).
\label{eq:feqo-double-absorption-statistics}
\end{equation}
这个式子的适用还要求相关高阶光子数矩有限，且 \(|\xi|\sqrt n\) 在初态有显著概率的 \(n\) 范围内小。固定平均数 \(\bar n=1\) 时，单光子数态给 \(\langle N(N-1)\rangle=0\)，相干态给 1，热态给 2；因此三者的一阶边带可以相同，而第二吸收边带依次为 \(0,|\xi|^4/4,|\xi|^4/2\) 的领先阶。若电子开始就占有多个边带、反冲使两次吸收失谐不同，须重新计算有序时间积分，不能继续把 \(\xi^2/2\) 当成普适振幅。

相同的 \(\bar n\) 在这个最低非零阶给相同人口，因而只比较弱场谱峰不能区分光统计。电子相干项却已经不同。按式~\eqref{eq:feqo-partial-trace-components} 对光求迹，
\begin{equation}
(\rho_e)_{1,0}=-\ii\xi\langle a\rangle+O(|\xi|^2),\qquad
(\rho_e)_{-1,0}=-\ii\xi^*\langle a^\dagger\rangle+O(|\xi|^2).
\label{eq:feqo-quantum-coherence-first-order}
\end{equation}
Fock 态 \(|n\rangle\) 与热态都满足 \(\langle a\rangle=0\)，所以一阶能量相干为零；相干态 \(|\alpha\rangle\) 满足 \(\langle a\rangle=\alpha\)，于是保留相位。比较时可选 \(n=1\)、\(|\alpha|^2=1\)、热平均数 \(\bar n=1\)，并固定同一个 \(\xi\)。人口先相同，电子相干先不同；更高阶人口还会读出 \(\langle N_\gamma(N_\gamma-1)\rangle\) 等统计矩。

对纯初始电子，保留式~\eqref{eq:feqo-weak-quantum-populations} 与式~\eqref{eq:feqo-quantum-coherence-first-order} 的最低非零阶，纯度为
\[
\operatorname{Tr}\rho_e^2
=1-2|\xi|^2\bigl(2\bar n+1-2|\langle a\rangle|^2\bigr)+O(|\xi|^3).
\]
Fock 和热态的括号是 \(2\bar n+1\)，相干态则是 1。该式只在扰动区内比较纯度；高阶项和具体脉冲相位可能改变更细的差别。

相干态为何会恢复上一章的经典场，可用位移算符看清。令 \(D(\alpha)=e^{\alpha a^\dagger-\alpha^*a}\)，则 \(D^\dagger aD=a+\alpha\)。在位移后的表象中，式~\eqref{eq:feqo-quantum-pinem-hamiltonian} 变为确定驱动 \(\hbar g\alpha B^\dagger+\mathrm{h.c.}\) 加上量子涨落耦合 \(\hbar gB^\dagger a+\mathrm{h.c.}\)。令 \(|\alpha|\to\infty\)、\(g\to0\)，同时固定 \(g\alpha\) 和作用时间，确定驱动不变而相对涨落尺度 \(1/|\alpha|\to0\)。此时电子演化趋向经典场的相位调制；仅说“光子数大”而不说明固定什么量，不构成一个受控极限。

式~\eqref{eq:feqo-joint-displacement} 还允许全阶比较三种光输入，而不在 \(|\xi|\sqrt{\bar n+1}\) 处截断。固定交换总量 \(N_\gamma+L_e=n\)，以末态光子数 \(m\geq0\) 标记联合基 \(|n-m,m\rangle\)。在这条链上，\(A\) 把 \(m\) 降一，\(A^\dagger\) 把 \(m\) 升一，作用系数与普通谐振子完全相同。定义位移矩阵元 \(D_{mn}(\zeta):=\langle m|e^{\zeta a^\dagger-\zeta^*a}|n\rangle\)。由式~\eqref{eq:feqo-joint-displacement} 的两个指数分别展开，先降 \(n-k\) 次、再升 \(m-k\) 次，得到有限和
\begin{equation}
D_{mn}(\zeta)=e^{-\mu/2}\sqrt{m!n!}
\sum_{k=0}^{\min(m,n)}
\frac{\zeta^{m-k}(-\zeta^*)^{n-k}}
{k!(m-k)!(n-k)!},\qquad \mu=|\zeta|^2.
\label{eq:feqo-displacement-matrix}
\end{equation}
这一式不要求读者事先认识关联 Laguerre 多项式。对入射 \(|0,n\rangle\)，末态振幅恰为 \(D_{mn}|n-m,m\rangle\)，故电子第 \(\ell\) 条边带的概率是 \(P_\ell^{(n)}=|D_{n-\ell,n}|^2\)（\(n-\ell\geq0\)）；其它 \(\ell\) 的概率为零。不同 \(\ell\) 留下不同光子数，所以电子约化态在能量基对角。

单光子 \(n=1\) 的前几项可以完全写开：
\[
P_{+1}^{(1)}=\mu e^{-\mu},\qquad
P_0^{(1)}=(1-\mu)^2e^{-\mu},\qquad
P_{-1}^{(1)}=\frac{\mu(2-\mu)^2}{2}e^{-\mu}.
\]
\(\mu=1\) 时中心边带恰被相消抽空，但概率由其它所有边带承接；只保留这三项不能强制它们的和等于一。若初态是真空 \(n=0\)，式~\eqref{eq:feqo-displacement-matrix} 退化为式~\eqref{eq:feqo-vacuum-joint-state} 的泊松分布。这个可直接检验的极限也固定了所有阶乘和符号。

相干光初态 \(|\alpha\rangle\) 的数态系数为 \(c_n=e^{-\bar n/2}\alpha^n/\sqrt{n!}\)，其中 \(\bar n=|\alpha|^2\)。同一末态光子数 \(m\) 可由不同初始 \(n=m+\ell\) 到达，所以联合振幅和电子约化矩阵分别是
\begin{equation}
C_{\ell m}=c_{m+\ell}D_{m,m+\ell}(\zeta),\qquad
(\rho_e)_{\ell\ell'}=\sum_{m\geq\max(0,-\ell,-\ell')}
C_{\ell m}C_{\ell'm}^* .
\label{eq:feqo-coherent-reduced-exact}
\end{equation}
负指标的 \(c_n\) 定义为零。这个和式保留来自不同初始光子数的相位干涉。热光初态则是对角混合，权重 \(p_n^{\rm th}=\bar n^n/(1+\bar n)^{n+1}\)，因此
\(P_\ell^{\rm th}=\sum_{n\geq\max(0,\ell)}p_n^{\rm th}|D_{n-\ell,n}|^2\)，电子非对角元为零。固定 \(\bar n=1\) 后，Fock 态 \(n=1\)、相干态 \(|\alpha|=1\)、热态可在完全相同的 \(\zeta\) 下比较，而不是把平均输入能量的差别误认为光统计。

这些无穷和也有透明的数值误差。若只保留初始光子数 \(n\leq N\)，任一电子人口的遗漏不超过输入光分布的尾概率 \(\sum_{n>N}p_n\)，因为每个固定 \(n\) 的条件概率不超过一。\(\bar n=1\) 的热态尾和正好是 \(2^{-(N+1)}\)；相干态尾和是 \(1-e^{-1}\sum_{n=0}^N1/n!\)。在计算前选择 \(N\) 使这个数小于目标误差，便有真正受控的截断。

\begin{exercise}
由式~\eqref{eq:feqo-displacement-matrix} 求 \(D_{0,1},D_{1,1},D_{2,1}\)，再用单光子结果解释 \(P_0^{(1)}\) 在 \(\mu=1\) 消失的原因。
\end{exercise}


若还记录某个光学结果，求和规则就改变。读到结果的电子子样本称条件态；它与完全不读光的约化态预测不同的实验。

\section{不读光与选定光学结果是两种实验}
\label{chap:feqo-conditioning}

式~\eqref{eq:feqo-partial-trace-components} 对所有未记录的光末态求和。如果实验还读出光场，可以只保留特定结果对应的电子。先看投影测量：光场投影 \(P_y=P_y^2=P_y^\dagger\) 代表结果 \(y\)，未归一化的条件电子态为
\begin{equation}
\widetilde\rho_{e|y}=\operatorname{Tr}_\gamma
\bigl[(I_e\otimes P_y)\rho_{e\gamma}(I_e\otimes P_y)\bigr],\qquad
p_y=\operatorname{Tr}\widetilde\rho_{e|y},\quad
\rho_{e|y}=\widetilde\rho_{e|y}/p_y.
\label{eq:feqo-conditional-electron-state}
\end{equation}
\(p_y=0\) 的结果不会发生，也就没有对应的条件态。把所有互斥结果的未归一化态相加，恢复不读光的 \(\rho_e\)；归一化后再相加则没有意义，必须带各自概率权重。

令前节二维示例在 \(gt=\pi/4\) 时的联合态为 \((|0,1\rangle-\ii|1,0\rangle)/\sqrt2\)。不读光得到 \(\rho_e=(|0\rangle\langle0|+|1\rangle\langle1|)/2\)，纯度 \(1/2\)。若光探测器投影到 \(|+\rangle_\gamma=(|0\rangle+|1\rangle)/\sqrt2\)，得到该结果的概率为 \(1/2\)，归一化的电子条件态为 \((|0\rangle-\ii|1\rangle)/\sqrt2\)，纯度为一。另一正交光学结果给相反的电子相位；忽略结果标签时，两种相位相加又消去。这就是“光场带走哪条路径信息”的可计算版本。

不必要求光学路径态恰好正交，才能量化这件事。考虑归一化纯态
\[
|\Psi\rangle=\sqrt p\,|0\rangle_e|u\rangle_\gamma
+e^{\ii\phi}\sqrt{1-p}\,|1\rangle_e|v\rangle_\gamma,
\qquad 0\leq p\leq1,
\]
其中 \(\langle u|u\rangle=\langle v|v\rangle=1\)，但 \(s:=\langle v|u\rangle\) 可以在零与单位模之间。对光求偏迹得到 \((\rho_e)_{01}=\sqrt{p(1-p)}e^{-\ii\phi}s\)，并由 \(2\times2\) 矩阵直接相乘求得
\begin{equation}
\operatorname{Tr}\rho_e^2=p^2+(1-p)^2+2p(1-p)|s|^2.
\label{eq:feqo-which-path-overlap}
\end{equation}
当 \(|u\rangle=|v\rangle\) 时 \(|s|=1\)，光没有记下电子走哪条能量路径，电子仍纯；当二者正交时 \(s=0\)，路径标签可由光完全区分，电子相干元消失。\(p\) 接近零或一时纯度也接近一，但那是因为电子几乎只走一条路径，不能误说环境“没有取得信息”。这个公式把路径信息、电子相干和纯度之间的关系落实到同一个可计算的内积。

光测量未必是正交投影。若结果 \(y\) 由作用于光场的算符 \(K_y\) 表示，则正算符 \(E_y=K_y^\dagger K_y\) 决定概率，\(\sum_yE_y=I\) 保证所有结果概率之和为一。把式~\eqref{eq:feqo-conditional-electron-state} 中的两个 \(P_y\) 分别换成 \(K_y\) 与 \(K_y^\dagger\) 就得到条件更新。结果效应 \(E_y\) 只决定概率；若要预测测量后继续演化的态，还必须给出具体 \(K_y\)。

电子谱变宽本身不证明纠缠：随机强度的经典光束也能混合出宽谱。这里的纠缠结论使用了已知的联合纯态及其约化纯度；实验验证则需要足以排除相应经典混合模型的电子—光相关或条件测量数据。

\begin{exercise}
对上面的联合态，把光投影改为 \(|-\rangle=(|0\rangle-|1\rangle)/\sqrt2\)，求电子条件态，并以两个结果的概率加权恢复无条件密度算符。
\end{exercise}

